% ======================================================================
%  Abstract
% ======================================================================
\begin{center}
  {\Large\bfseries Abstract}
\end{center}
\vspace{0.4cm}

\noindent
At the sciNeurotech Lab, neural responses to stimulation parameters are unknown
and expensive to evaluate, which makes Gaussian-process Bayesian optimization
(GPBO) a natural tool for driving stimulation experiments. Crucially, the
\emph{same} biological response can later serve very different scientific goals:
matching a prescribed target level, trading off several competing responses,
optimizing under constraints, or exploiting prior knowledge about promising
regions. This raises a modelling question that cuts across all these tasks:
\emph{where should task-specific information enter the Bayesian-optimization
pipeline?}

This report develops and studies a \emph{separation-of-concerns} (SoC) principle:
model the underlying phenomenon $f$ faithfully with a Gaussian process, and
encode the task-specific transformation at the decision (acquisition) level
whenever possible, rather than folding it into the statistical model. We contrast
this \emph{separated} formulation with the \emph{task-in-model} formulation, in
which a GP is fitted directly to the transformed objective $\prm\circ f$, and we
ask when moving task information into the model changes the assumptions on which
GPBO relies.

On the theoretical side, the central contribution is a study of
\emph{task-induced misspecification} for target matching: the transformation
$g=\abs{f-T}$ creates a cusp at the target level that, under smooth-kernel
assumptions, pushes $g$ outside the assumed reproducing-kernel Hilbert space
(RKHS). Because mere non-membership is a weak statement, we sharpen it into a
\emph{quantitative approximation barrier} that lower-bounds how well $g$ can be
approximated by an RKHS ball of controlled radius. We complement this with a
positive result: keeping the GP on $f$ and pushing its confidence interval
through the task recovers a GP-UCB-type regret guarantee. The same
pushforward-confidence viewpoint is then extended to multi-objective
scalarization and to prior-guided optimization, where we analyse the gap between
the (inaccessible) global RKHS criterion and coefficients estimated only from
adaptively sampled data, and exhibit a self-confirming feedback mechanism as well as a possible convergence failure when the prior knowledge is put directly in the model.

On the experimental side, three controlled studies (target matching,
multi-objective optimization, prior-guided optimization) compare the separated
and task-in-model formulations on standard benchmarks, with matched budgets,
shared initial designs and paired statistical tests. They consistently favour
separation, most sharply on deceptive low-dimensional problems, while narrowing
in higher dimensions.

The main limitation is deliberately stated: our results establish
\emph{misspecification} and \emph{approximation} barriers, not a general lower
regret bound for the task-in-model strategy. RKHS misspecification does not by
itself imply empirical failure, and a true lower bound for task-in-model GP-UCB
remains open.

\vspace{0.4cm}
\noindent\textbf{Keywords:} Bayesian optimization, Gaussian processes, RKHS,
model misspecification, target matching, multi-objective optimization,
prior-guided optimization, neurostimulation.
